\section{Introduction}\label{sec:intro}

Camera-equipped smart glasses are a fast-growing device category~\cite{idc2026xr,counterpoint2026aiglasses}, and an assistant on them is most useful when the wearer can refer to things in view. Pointing is the most direct way to do so: ``what is \emph{that}?'', ``remind me to buy \emph{this}''. Systems have combined speech with pointing since Bolt's \emph{Put-that-there}~\cite{bolt1980put}, and recent wearable assistants use pointing and gaze to resolve such pronouns~\cite{lee2024gazepointar}. Before any of this works on glasses, the device must answer a spatial question from its own camera: \emph{which object is the wearer pointing at?}

Three things make this hard on glasses. First, the camera sees the wearer's hand from behind and often only in part, and the target may overlap the fingertip in the image or lie far beyond it. Second, other people's hands can be in view and must not be mistaken for the wearer's. Third, the budget is small: commercial smart glasses carry batteries of a few hundred milliampere-hours that the camera, audio and radio share~\cite{yuan2026aiglasses}, while sending frames to a cloud model adds network latency and moves imagery of bystanders off the device.

Pointing-target grounding in egocentric images is an active research area, for which EgoPoint-Ground provides a large benchmark~\cite{egopoint}. The most accurate methods are large: multimodal LLMs that reason step by step from the hand to a ray and then to the target~\cite{egopoint,li2026pointvgr}, and a 226M-parameter transformer that models the pointing ray explicitly~\cite{li2026vistaref}. Some benchmark tasks even give the model the ground-truth hand box~\cite{egopoint}. These systems reach high accuracy on a GPU, but none is designed for glasses-class hardware, and to our knowledge none evaluates other people's hands in view. Egocentric fingertip systems, including one that runs on a phone, stop at the fingertip~\cite{huang2016pointing,jain2019gestarlite}, and tiny detectors that run on smart glasses find generic objects~\cite{moosmann2024tinyissimoyolo} but do not decide which one a gesture refers to.

We therefore ask how much of this spatial step a model small enough for the glasses can do. SpatialNet takes one RGB frame and nothing else: no text prompt, no hand box, no depth, and no external hand detector or pose estimator. In one forward pass of a compact convolutional network with no self-attention and no token-by-token decoding (0.72M parameters, 0.23 GFLOPs, 1.6\,MB in FP16), it predicts the wearer's hand, the fingertip, the pointing direction and the box of the pointed-at object, together with a hand-presence gate. Its design also has pointing and target gates, but EgoPoint provides no negative examples to train or test them. The model is at least 300 times smaller than VistaRef and about 10{,}000 times smaller than the 7B multimodal LLMs.

One design choice targets precision at this size: \emph{instance-anchored heads}. The heads form a chain in which each one reads a shared feature map at the location the previous one predicted, so that the fingertip, direction and target belong to the same hand. We also give the target head the geometry of the predicted ray, through a zero-initialised branch, and then test whether a model this small uses it rather than assuming so. Training pastes unlabelled hands of other people into the frames and penalises responses to them, so that the model learns to prefer the wearer's hand without labelled bystander hands.

We train and test on EgoPoint and report every target metric by the gap between the fingertip and the target in the image, because in about 70\% of images the fingertip lies inside the target box. We compare SpatialNet with published large models, with simple fingertip-based references, and with compact detectors (YOLO26n and NanoDet-Plus-m) that we trained on the same split at the same compute and scored with the same code. Because at least 41\% of the EgoPoint test images are COCO photographs, we separate detectors pretrained on COCO from those that, like SpatialNet, never saw COCO data, and test whether the overlap matters. We ablate each component, test behaviour when other people's hands are pasted into the frames, and measure latency, memory, battery drain and accuracy on the camera frames of a pre-production pair of smart glasses. We do not claim a more accurate localiser than large models, or tighter boxes than a compact detector. We show what a sub-million-parameter, image-only model can do on smart-glasses hardware, where it fails, and which design choices matter at that scale.

\paragraph{Contributions.}
\begin{enumerate}
\item \emph{A deployable, image-only pointing-target localiser for smart glasses.} A 0.72M-parameter CNN returns the wearer's hand, fingertip, direction and target box in one pass, and runs in 60--80\,ms per frame on a pre-production glasses unit, where we also check its accuracy on the unit's own camera frames (Sections~\ref{sec:method}, \ref{sec:device}, \ref{sec:ondevice} and~\ref{sec:realframes}).
%% If the real-frame check (Section 7.8) cannot be done in time, delete the clause ", where we also check ...", and Section 7.8.
\item \emph{A controlled characterisation of accuracy against cost at wearable scale.} Beyond published large models, we compare with compact detectors trained under our protocol, including a pretraining control, and report an end-to-end score, a resolution study, results by fingertip--target gap, retraining variation and a duplicate audit. At the same compute and with the same ImageNet pretraining, a detector draws tighter boxes (63.2\% against 58.9\% P@0.5), but SpatialNet places its target centre on the labelled object more often (84.8\% against 81.6--82.2\% for four detector trainings, all paired intervals excluding zero), and in a downstream test an off-the-shelf recogniser names the target from SpatialNet's crops as often as from the detectors', except when the target lies far from the fingertip. We also show that at least 41\% of the EgoPoint test images are COCO train2014 photographs, which any benchmark on this dataset should report, although COCO pretraining gains nothing extra on them (Sections~\ref{sec:data} and \ref{sec:main}--\ref{sec:heads-results}).
\item \emph{Mechanism-level findings for small pointing models.} Input resolution is the largest lever, at no parameter cost. Instance anchoring makes the hand box and the direction more precise than a plain model of the same size, without changing coarse accuracy. Distractor pasting keeps accuracy under synthetic interference at no cost on clean images, and the recipe transfers unchanged to a generic detector. An explicit ray prior learns mainly fingertip proximity when most targets overlap the fingertip; neither constraining it to point forward, nor up-weighting large-gap targets, nor widening the size sample fixes large-gap or box-extent errors (Sections~\ref{sec:detectors} and \ref{sec:resolution}--\ref{sec:own}).
\end{enumerate}
